\documentclass[11pt,a4paper]{article}

\usepackage{thga-db}
\usepackage{listings}

\renewcommand{\thgaKurs}{Databases}
\renewcommand{\thgaDatum}{Summer Term 2026}

\begin{document}
\thispagestyle{empty}
\thgaTitel{Term Project – Documentation }%
{Wartungsmangement Datenbankanwendung}%

\begin{infobox}
\textbf{Author:} Niklas Stratenwerth  \\
\textbf{Module:} Introduction to Database Management Systems · THGA Bochum \\
\textbf{Stack:} PostgreSQL · FastAPI · Docker Compose · tkinter (\texttt{.deb}) \\
\textbf{Repository:} \url{https://github.com/Niklas0511/DBMS_Project} \\
\end{infobox}

\tableofcontents
\vspace{1em}

% ============================================================
\section{Introduction and motivation}
% ============================================================

Ob Geräte gewartet werden wird häufig nicht dokumentiert, oder ob diese ersetzt wurden.
Desweiteren ist häufig unklar, welche Geräte überhaupt an einem Standort vorhanden sind.
Wenn dies dolumentiert wird ist es meist Formlos und von Standort zu Standort unterschiedlich.
Die Anwendung soll die Dokumentation der Wartung von Geräten vereinfachen und standardisieren.

\textbf{Wartungsmanagement} Hält sowohl die Geräte als auch die Wartungsdaten in einer zentralen Datenbank fest.
Es verwaltet \textbf{Geräte}, \textbf{Wartungen} and \textbf{Standorte}, Die Daten können durch eine Desktop Anwendung abgerufen
und modifiziert werden.

\begin{beispielbox}
\textbf{Anwendungsszenario.} Ein Techniker möchte sehen, welche Wartungen aktuell geplant sind.
Dazu kann er in der Desktio Anwendung nachschauen, welche Wartungen für diesen Monat geplant sind.
Er kann auch sehen, welche Geräte an einem Standort vorhanden sind und welche Wartungen bereits durchgeführt wurden. 
\end{beispielbox}

% ============================================================
\section{Data model}
% ============================================================

Die Domäne umfasst die Entitätstypen \textbf{Standort},
\textbf{Techniker}, \textbf{Gerät}, \textbf{Wartungsplan}
und \textbf{Wartung} sowie zwei Zuordnungstabellen.

Ein Standort kann mehrere Geräte beherbergen, während jedes Gerät
genau einem Standort zugeordnet ist. Zwischen Standorten und
Technikern besteht eine N:M-Beziehung, die durch die Tabelle
\texttt{standort\_techniker} aufgelöst wird.

Die Tabelle \texttt{typ\_techniker} ordnet Technikern die Gerätetypen
zu, für die sie zuständig sind. Ein Techniker kann für mehrere
Gerätetypen zuständig sein, und ein Gerätetyp kann mehreren
Technikern zugeordnet werden. Gerätetypen werden als Zeichenketten
gespeichert; eine eigene Entität für Gerätetypen existiert nicht.

Ein Wartungsplan beschreibt eine geplante Wartung für genau ein
Gerät durch genau einen Techniker. Eine Wartung dokumentiert
hingegen eine tatsächlich durchgeführte Arbeit und verweist
ebenfalls auf genau ein Gerät und einen Techniker.
Sie kann optional einem Wartungsplan zugeordnet werden.

\begin{beispielbox}
\textbf{ER-Modell (textuell).}

\begin{itemize}[leftmargin=1.4em, topsep=2pt]
  \item \textbf{Standort} (1) -- beherbergt -- (N)
        \textbf{Gerät}

  \item \textbf{Standort} (1) -- ist zugeordnet -- (N)
        \textbf{Standort-Techniker-Zuordnung}

  \item \textbf{Techniker} (1) -- ist zugeordnet -- (N)
        \textbf{Standort-Techniker-Zuordnung}

  \item \textbf{Techniker} (1) -- ist zugeordnet -- (N)
        \textbf{Gerätetyp-Techniker-Zuordnung}

  \item \textbf{Gerät} (1) -- hat -- (N)
        \textbf{Wartungsplan}

  \item \textbf{Techniker} (1) -- ist eingeteilt für -- (N)
        \textbf{Wartungsplan}

  \item \textbf{Gerät} (1) -- erhält -- (N)
        \textbf{Wartung}

  \item \textbf{Techniker} (1) -- führte durch -- (N)
        \textbf{Wartung}

  \item \textbf{Wartungsplan} (0..1) -- ist zugeordnet zu -- (N)
        \textbf{Wartung}
\end{itemize}

Damit besteht zwischen \textbf{Standort} und \textbf{Techniker}
eine N:M-Beziehung, aufgelöst durch
\texttt{standort\_techniker}.

Jede Wartung verweist auf höchstens einen Wartungsplan.
Das vorliegende Schema erlaubt jedoch mehrere Wartungen pro
Wartungsplan, da \texttt{wartungsplan\_id} in der Tabelle
\texttt{wartung} nicht eindeutig sein muss.
\end{beispielbox}

\subsection{Relationales Schema}

Aus dem ER-Modell ergeben sich die folgenden Relationen.
Primärschlüssel sind unterstrichen, Fremdschlüssel kursiv dargestellt.
Attribute, die zu beiden Schlüsselarten gehören, sind sowohl
unterstrichen als auch kursiv gesetzt.

\begin{itemize}[leftmargin=1.4em]
  \item \texttt{standort}(\underline{id}, plz, stadt,
        strasse, hausnummer)

  \item \texttt{techniker}(\underline{id}, name)

  \item \texttt{geraet}(\underline{seriennummer}, hersteller,
        typ, wartungsintervall, \emph{standort\_id})

  \item \texttt{standort\_techniker}(
        \underline{\emph{standort\_id}},
        \underline{\emph{techniker\_id}})

  \item \texttt{typ\_techniker}(
        \underline{geraetetyp},
        \underline{\emph{techniker\_id}})

  \item \texttt{wartungsplan}(\underline{id}, datum,
        \emph{geraet\_seriennummer}, \emph{techniker\_id})

  \item \texttt{wartung}(\underline{id}, datum,
        ersatzgeraet, \emph{geraet\_seriennummer},
        \emph{techniker\_id}, \emph{wartungsplan\_id})
\end{itemize}

\textbf{Integritätsbedingungen.}
Die IDs von Standorten, Technikern, Wartungsplänen und Wartungen
werden automatisch erzeugt. Geräte werden über ihre Seriennummer
identifiziert. Die zusammengesetzten Primärschlüssel der
Zuordnungstabellen verhindern doppelte Zuordnungen.

Das Wartungsintervall eines Geräts muss größer als null sein.
Die Kombination aus Datum und Geräteseriennummer ist in
\texttt{wartungsplan} eindeutig. Für ein Gerät kann deshalb
höchstens eine Wartung pro Datum geplant werden.

Die Attribute \texttt{ersatzgeraet} und
\texttt{wartungsplan\_id} dürfen in \texttt{wartung}
leer sein. \texttt{ersatzgeraet} ist ein optionales Textfeld
und kein Fremdschlüssel auf die Gerätetabelle.
Auch \texttt{geraetetyp} in \texttt{typ\_techniker}
ist kein Fremdschlüssel.

\textbf{Referenzielle Integrität.}
Fremdschlüssel sichern die Existenz referenzierter Datensätze.
Änderungen referenzierter Schlüssel werden durch
\texttt{ON UPDATE CASCADE} weitergegeben.

Ein Standort mit zugeordneten Geräten kann nicht gelöscht werden.
Ebenso verhindern vorhandene Wartungen das Löschen des betreffenden
Geräts oder Technikers. Ein Techniker kann außerdem nicht gelöscht
werden, solange ihm Wartungspläne zugeordnet sind.

Beim Löschen eines Geräts werden dessen Wartungspläne gelöscht,
sofern keine andere Bedingung das Löschen des Geräts verhindert.
Beim Löschen eines Wartungsplans bleiben durchgeführte Wartungen
erhalten; ihre Planreferenz wird auf \texttt{NULL} gesetzt.
Zuordnungen in den beiden Zuordnungstabellen werden beim Löschen
des jeweils referenzierten Standorts oder Technikers automatisch
entfernt.

\subsection{Normalisierung}

Alle Attribute enthalten einzelne Werte; mehrfache Zuständigkeiten
werden durch separate Zuordnungszeilen dargestellt.
Damit erfüllt das Schema die \textbf{erste Normalform}.

Die Entitätstabellen besitzen jeweils einen einspaltigen
Primärschlüssel. Die beiden Zuordnungstabellen bestehen
ausschließlich aus den Attributen ihres zusammengesetzten
Primärschlüssels. Partielle Abhängigkeiten von Nichtschlüsselattributen
liegen daher nicht vor; die \textbf{zweite Normalform} ist erfüllt.

Standortdaten werden ausschließlich in \texttt{standort},
Technikernamen in \texttt{techniker} und Gerätestammdaten
in \texttt{geraet} gespeichert. Wartungspläne und Wartungen
referenzieren diese Datensätze über ihre Schlüssel.

Unter der Annahme, dass keine weiteren fachlichen funktionalen
Abhängigkeiten bestehen, erfüllt das Schema die
\textbf{dritte Normalform}. 

Gerät und Techniker einer durchgeführten Wartung werden unabhängig
von einer optionalen Planreferenz gespeichert. Das Schema erzwingt
nicht, dass diese Angaben mit den Angaben des referenzierten
Wartungsplans übereinstimmen.
(L 04).

% ============================================================
\section{Database (PostgreSQL)}
% ============================================================

Das Schema wird durch einen Init-Script erstellt, der von PostgreSQL bei der ersten Start ausgeführt wird (L 07, L 09).

\begin{lstlisting}[language=SQL, caption={init.sql (excerpt) -- schema and seed data}]
CREATE TABLE standort (
    id BIGINT GENERATED ALWAYS AS IDENTITY PRIMARY KEY,
    plz VARCHAR(10) NOT NULL,
    stadt VARCHAR(100) NOT NULL,
    strasse VARCHAR(100) NOT NULL,
    hausnummer VARCHAR(20) NOT NULL
);

CREATE TABLE techniker (
    id BIGINT GENERATED ALWAYS AS IDENTITY PRIMARY KEY,
    name VARCHAR(150) NOT NULL
);

CREATE TABLE geraet (
    seriennummer VARCHAR(100) PRIMARY KEY,
    hersteller VARCHAR(100) NOT NULL,
    typ VARCHAR(100) NOT NULL,
    wartungsintervall INTEGER NOT NULL CHECK (wartungsintervall > 0),
    standort_id BIGINT NOT NULL REFERENCES standort(id)
        ON UPDATE CASCADE ON DELETE RESTRICT
);

CREATE TABLE standort_techniker (
    standort_id BIGINT NOT NULL REFERENCES standort(id)
        ON UPDATE CASCADE ON DELETE CASCADE,
    techniker_id BIGINT NOT NULL REFERENCES techniker(id)
        ON UPDATE CASCADE ON DELETE CASCADE,
    PRIMARY KEY (standort_id, techniker_id)
);

CREATE TABLE typ_techniker (
    geraetetyp VARCHAR(100) NOT NULL,
    techniker_id BIGINT NOT NULL REFERENCES techniker(id)
        ON UPDATE CASCADE ON DELETE CASCADE,
    PRIMARY KEY (geraetetyp, techniker_id)
);

CREATE TABLE wartungsplan (
    id BIGINT GENERATED ALWAYS AS IDENTITY PRIMARY KEY,
    datum DATE NOT NULL,
    geraet_seriennummer VARCHAR(100) NOT NULL REFERENCES geraet(seriennummer)
        ON UPDATE CASCADE ON DELETE CASCADE,
    techniker_id BIGINT NOT NULL REFERENCES techniker(id)
        ON UPDATE CASCADE ON DELETE RESTRICT,
    UNIQUE (datum, geraet_seriennummer)
);

CREATE TABLE wartung (
    id BIGINT GENERATED ALWAYS AS IDENTITY PRIMARY KEY,
    datum DATE NOT NULL,
    ersatzgeraet VARCHAR(100),
    geraet_seriennummer VARCHAR(100) NOT NULL REFERENCES geraet(seriennummer)
        ON UPDATE CASCADE ON DELETE RESTRICT,
    techniker_id BIGINT NOT NULL REFERENCES techniker(id)
        ON UPDATE CASCADE ON DELETE RESTRICT,
    wartungsplan_id BIGINT REFERENCES wartungsplan(id)
        ON UPDATE CASCADE ON DELETE SET NULL
);

CREATE INDEX idx_geraet_standort ON geraet (standort_id);
CREATE INDEX idx_wartungsplan_datum ON wartungsplan (datum);
CREATE INDEX idx_wartung_geraet ON wartung (geraet_seriennummer);

\end{lstlisting}

% ============================================================
\section{REST API (FastAPI)}
% ============================================================

Der Zugriff auf die Datenbank erfolgt nur über die REST API (L 08). Read endpoints
sind öffentlich; write endpoints brauchen den header \texttt{X-API-Key} (cf.\ DBMS\_09).

\renewcommand{\arraystretch}{1.2}
\begingroup
\small
\renewcommand{\arraystretch}{1.2}

\begin{tabularx}{\linewidth}{
    @{}
    l
    >{\raggedright\arraybackslash}p{5.8cm}
    >{\raggedright\arraybackslash}X
    c
    @{}
}
\toprule
\textbf{Methode} & \textbf{Pfad} & \textbf{Zweck} & \textbf{Key} \\
\midrule

\texttt{GET}
& \nolinkurl{/health}
& Datenbankverbindung prüfen
& -- \\

\texttt{GET}
& \nolinkurl{/techniker/{techniker_id}/wartungsplan}
& Wartungspläne eines Technikers mit Geräte- und Standortdaten lesen (JOIN)
& -- \\

\texttt{GET}
& \nolinkurl{/geraete/{seriennummer}/wartungen}
& Wartungshistorie eines Geräts mit Technikerdaten lesen (JOIN)
& -- \\

\texttt{GET}
& \nolinkurl{/geraete/{seriennummer}}
& Gerät mit Standort und nächster offener Wartung ab heute lesen (JOIN, Unterabfrage)
& -- \\

\texttt{GET}
& \nolinkurl{/standorte}
& Alle Standorte auflisten
& -- \\

\texttt{GET}
& \nolinkurl{/standorte/{standort_id}/geraete}
& Geräte eines Standorts auflisten
& -- \\

\texttt{GET}
& \nolinkurl{/wartungsplan}
& Wartungspläne im Zeitraum mit Geräte-, Techniker-, Standortdaten und Durchführungsstatus lesen (JOIN, EXISTS)
& -- \\

\texttt{GET}
& \nolinkurl{/wartungsplan/ueberfaellig}
& Überfällige, offene Wartungspläne mit Tagen der Überschreitung auflisten (JOIN, NOT EXISTS)
& -- \\

\texttt{GET}
& \nolinkurl{/techniker}
& Techniker mit zuständigen Standorten und Gerätetypen auflisten (Unterabfragen)
& -- \\

\texttt{POST}
& \nolinkurl{/wartungen}
& Durchgeführte Wartung erfassen, optional einem Wartungsplan zuordnen
& Ja \\

\texttt{POST}
& \nolinkurl{/geraete}
& Neues Gerät mit Standortzuordnung anlegen
& Ja \\

\texttt{POST}
& \nolinkurl{/wartungsplan}
& Wartung für ein Gerät mit Datum und Techniker planen
& Ja \\

\texttt{POST}
& \nolinkurl{/standorte}
& Neuen Standort anlegen
& Ja \\

\texttt{POST}
& \nolinkurl{/techniker}
& Techniker mit Standort- und Gerätetypzuständigkeiten anlegen
& Ja \\

\texttt{PUT}
& \nolinkurl{/techniker/{techniker_id}}
& Name und Zuständigkeiten eines Technikers ersetzen
& Ja \\

\texttt{PUT}
& \nolinkurl{/wartungsplan/{wartungsplan_id}}
& Datum und Techniker eines Wartungsplans ändern
& Ja \\

\texttt{PATCH}
& \nolinkurl{/geraete/{seriennummer}}
& Einzelne Geräteeigenschaften ändern
& Ja \\

\texttt{DELETE}
& \nolinkurl{/wartungsplan/{wartungsplan_id}}
& Wartungsplan löschen
& Ja \\

\texttt{GET}
& \nolinkurl{/statistik/wartungen-pro-techniker}
& Durchgeführte Wartungen je Techniker im Zeitraum zählen (LEFT JOIN, Aggregation)
& -- \\

\texttt{GET}
& \nolinkurl{/statistik/wartungen-pro-standort}
& Durchgeführte Wartungen und offene Wartungspläne je Standort im Zeitraum zählen (LEFT JOIN, Aggregation)
& -- \\

\bottomrule
\end{tabularx}

\medskip
\noindent
Die Endpunkte \nolinkurl{/wartungsplan},
\nolinkurl{/statistik/wartungen-pro-techniker} und
\nolinkurl{/statistik/wartungen-pro-standort}
erwarten die Query-Parameter \texttt{von} und \texttt{bis}
im Format \texttt{YYYY-MM-DD}.
Die Grenzen des Zeitraums sind eingeschlossen.

\endgroup


Alle sachreibenden Endpunkte sind durch einen API-Key-Guard geschützt -- adapted
from the DBMS\_09 exercise:

\begin{lstlisting}[language=Python, caption={API-key guard (adapted from DBMS\_09)}]
from fastapi import Header, HTTPException
import os

def require_api_key(
    x_api_key: str | None = Header(default=None, alias="X-API-Key"),
) -> None:
    if x_api_key is None or not secrets.compare_digest(x_api_key, API_KEY):
        raise HTTPException(
            status_code=status.HTTP_401_UNAUTHORIZED,
            detail="Ungueltiger oder fehlender API-Key",
        )
\end{lstlisting}

% ============================================================
\section{Frontend (tkinter, packaged as .deb)}
% ============================================================

Das Frontend ist eine tkinter-Anwendung, welche über die API auf die Datenbank zugreift.
Hierbei kann der API Key in der \texttt{.env} Datei hinterlegt werden, oder bei Start der Anwendung in
dem X-API-Key-Feld eingegeben werden.

\textbf{Ausführen} Mittels uv ist es möglich, die Anwendung zu starten, ohne dass die Abhängigkeiten installiert sein müssen.
Der genaue Ablauf ist in der README Datei im Repository beschrieben.

\textbf{Packaging.}
Das Frontend kann als Debian-Paket bereitgestellt und auf einem
kompatiblen Ubuntu- oder Debian-System mit grafischer Oberfläche
installiert werden. PyInstaller bündelt die Python-Anwendung und
ihre Python-Abhängigkeiten. Anschließend erstellt
\texttt{dpkg-deb} daraus das installierbare \texttt{.deb}-Paket.

Voraussetzung sind die installierten Projektabhängigkeiten sowie
PyInstaller. Die Paketstruktur im Verzeichnis \texttt{deb-paket}
muss einschließlich Paketinformationen, Startskript und
Menüeintrag vorbereitet sein. Die folgenden Befehle werden
im Verzeichnis \texttt{FE} ausgeführt.

\begin{lstlisting}[caption={Erstellen und Installieren des Debian-Pakets}]
uv add --dev pyinstaller

uv run python -m PyInstaller --noconfirm --clean --onedir \
  --name wartungsmanagement main.py

mkdir -p deb-paket/opt/wartungsmanagement
cp -a dist/wartungsmanagement/. \
  deb-paket/opt/wartungsmanagement/

dpkg-deb --build --root-owner-group \
  deb-paket wartungsmanagement_1.0.0.deb

sudo apt install ./wartungsmanagement_1.0.0.deb
\end{lstlisting}



% ============================================================
\section{Operation and deployment}
% ============================================================
Das Backend läuft in einem Docker-Container. Alle Zugangsdaten sind in einer \texttt{.env} Datei gespeichert,
welche nicht in das Repository committet wird. Dafür wird aber eine \texttt{.env.example} Datei bereitgestellt, welche die Struktur der \texttt{.env} Datei zeigt.
Die genauen Schritte zum Starten der Anwendung sind in den entsprechenden README Datei im Repository beschrieben.

\begin{lstlisting}[caption={docker-compose.yml (excerpt, after L 09)}]
services:
  db:
    image: postgres:17-alpine
    restart: unless-stopped
    environment:
      POSTGRES_DB: ${POSTGRES_DB}
      POSTGRES_USER: ${POSTGRES_USER}
      POSTGRES_PASSWORD: ${POSTGRES_PASSWORD}
    volumes:
      - postgres_data:/var/lib/postgresql/data
      - ./database/init:/docker-entrypoint-initdb.d:ro
    healthcheck:
      test: ["CMD-SHELL", "pg_isready -U ${POSTGRES_USER} -d ${POSTGRES_DB}"]
      interval: 5s
      timeout: 5s
      retries: 10

  api:
    build:
      context: .
      dockerfile: Dockerfile
    restart: unless-stopped
    environment:
      DATABASE_URL: postgresql://${POSTGRES_USER}:${POSTGRES_PASSWORD}@db:5432/${POSTGRES_DB}
      API_KEY: ${API_KEY}
    ports:
      - "8000:8000"
    depends_on:
      db:
        condition: service_healthy

volumes:
  postgres_data:

\end{lstlisting}

\textbf{Deployment steps}
\begin{enumerate}[leftmargin=1.6em, topsep=2pt]
  \item \texttt{.env} erstellen und Zugangsdaten eintragen
  \item \texttt{docker compose up -d} -- starte Datenbank und API.
  \item \texttt{sudo dpkg -i wartungsmanagement\_1.0.0.deb} -- installiere das Frontend.
  \item Launch \texttt{wartungsmanagement}, Wenn nicht schon in der \texttt{.env} Datei hinterlegt, kann der
        \texttt{X-API-Key} eigegeben werden.
\end{enumerate}

% ============================================================
\section{Reflection}
% ============================================================

Die Aufteilung in Backend und Frontend hat den Vorteil, dass die Datenbank und die API auf einem Server laufen können, 
während das Frontend auf einem anderen Rechner läuft. Desweiteren gab es den Dateien auf natürliche Weise eine Trennung und 
damit dem Projektordner von Anfang an eine Struktur. Docker und uv ermöglichen die Ausführung auf verschhiedenen Systemen, was
mir zwar in der Entwicklung nicht geholfen hat, aber für die spätere Nutzung der Anwendung von Vorteil ist. In Zukunft könnte
man die Anwendung noch erweitern um z.B. eine Benutzerverwaltung, sodass nicht jeder Techniker die Datenbank einsehen kann, 
sondern nur die für ihn relevanten Daten oder eine Möglichkeit Wartungen teilautomatisiert zu planen.

% ============================================================
\section*{Sources and reuse}
% ============================================================
\addcontentsline{toc}{section}{Sources and reuse}

Reused with attribution, as permitted:
\begin{itemize}[leftmargin=1.4em, topsep=2pt]
  \item API-key guard and Compose file structure -- adapted from
        \textbf{DBMS\_09} and \textbf{lecture 09}.
  \item tkinter \texttt{Treeview} layout -- adapted from
        \textbf{DBMS\_09 (frontend)} and \textbf{lecture 10}.
  \item PyInstaller / \texttt{fpm} packaging commands -- after \textbf{lecture
        13}.
\end{itemize}

\end{document}
